\begin{proof}
(i) Como $A$ es un subanillo de $R$, $A_\mathfrak p$ es un subanillo
de $R_\mathfrak p$ para todo ideal primo $\mathfrak p$ de $A$, y
$R_\mathfrak p$ es un anillo de fracciones de $A_\mathfrak p$; además,
$R_\mathfrak p$ es necesariamente un producto directo de un número
finito de cuerpos, de modo que todo elemento que no sea divisor de
$0$ en $R_\mathfrak p$ es invertible; esto demuestra que
$R_\mathfrak p$ es el anillo total de fracciones de
$A_\mathfrak p$. Si $A$ está integralmente cerrado en $R$,
$A_\mathfrak p$ está, por tanto, integralmente cerrado en
$R_\mathfrak p$; pero, si $R_\mathfrak p$ es un producto directo de al
menos 2 cuerpos, la clausura integral de $A_\mathfrak p$ en
$R_\mathfrak p$ es un producto directo de al menos 2 anillos no
reducidos a $0$, lo cual es absurdo, puesto que $A_\mathfrak p$ es un
anillo local; por consiguiente, $R_\mathfrak p$ es necesariamente un
cuerpo y $A_\mathfrak p$ es un dominio entero integralmente cerrado,
lo que por definición significa que $A$ es normal. Recíprocamente,
si $A$ es normal, $R_\mathfrak p$, anillo total de fracciones del
dominio entero $A_\mathfrak p$, es un cuerpo y $A_\mathfrak p$ está
integralmente cerrado en $R_\mathfrak p$; si $x\in R$ es un elemento
entero sobre $A$, su imagen en cada $R_\mathfrak p$ es entera sobre
$A_\mathfrak p$ y, por tanto, pertenece a $A_\mathfrak p$; se concluye
que $x\in A$ (Bourbaki, \emph{Alg. comm.}, cap. II, § 3, n\textsuperscript{o} 3, cor. 1 del
teor. 1) y, en consecuencia, $A$ está integralmente cerrado en $R$.

(ii) Como $A\subset A_\lambda\subset R$ para todo $\lambda$, $R$ es
también el anillo total de fracciones de cada $A_\lambda$. Dada la
caracterización de los anillos normales que tienen a $R$ por anillo
total de fracciones, establecida en (i), la afirmación (ii) se deduce
de Bourbaki, \emph{Alg. comm.}, cap. V, § 1, n\textsuperscript{o} 3, prop. 12.
\end{proof}
